% Options for packages loaded elsewhere
\PassOptionsToPackage{unicode}{hyperref}
\PassOptionsToPackage{hyphens}{url}
\documentclass[
  11pt,
]{article}
\usepackage{xcolor}
\usepackage[margin=1in]{geometry}
\usepackage{amsmath,amssymb}
\setcounter{secnumdepth}{5}
\usepackage{iftex}
\ifPDFTeX
  \usepackage[T1]{fontenc}
  \usepackage[utf8]{inputenc}
  \usepackage{textcomp} % provide euro and other symbols
\else % if luatex or xetex
  \usepackage{unicode-math} % this also loads fontspec
  \defaultfontfeatures{Scale=MatchLowercase}
  \defaultfontfeatures[\rmfamily]{Ligatures=TeX,Scale=1}
\fi
\usepackage{lmodern}
\ifPDFTeX\else
  % xetex/luatex font selection
\fi
% Use upquote if available, for straight quotes in verbatim environments
\IfFileExists{upquote.sty}{\usepackage{upquote}}{}
\IfFileExists{microtype.sty}{% use microtype if available
  \usepackage[]{microtype}
  \UseMicrotypeSet[protrusion]{basicmath} % disable protrusion for tt fonts
}{}
\makeatletter
\@ifundefined{KOMAClassName}{% if non-KOMA class
  \IfFileExists{parskip.sty}{%
    \usepackage{parskip}
  }{% else
    \setlength{\parindent}{0pt}
    \setlength{\parskip}{6pt plus 2pt minus 1pt}}
}{% if KOMA class
  \KOMAoptions{parskip=half}}
\makeatother
\usepackage{color}
\usepackage{fancyvrb}
\newcommand{\VerbBar}{|}
\newcommand{\VERB}{\Verb[commandchars=\\\{\}]}
\DefineVerbatimEnvironment{Highlighting}{Verbatim}{commandchars=\\\{\}}
% Add ',fontsize=\small' for more characters per line
\usepackage{framed}
\definecolor{shadecolor}{RGB}{248,248,248}
\newenvironment{Shaded}{\begin{snugshade}}{\end{snugshade}}
\newcommand{\AlertTok}[1]{\textcolor[rgb]{0.94,0.16,0.16}{#1}}
\newcommand{\AnnotationTok}[1]{\textcolor[rgb]{0.56,0.35,0.01}{\textbf{\textit{#1}}}}
\newcommand{\AttributeTok}[1]{\textcolor[rgb]{0.13,0.29,0.53}{#1}}
\newcommand{\BaseNTok}[1]{\textcolor[rgb]{0.00,0.00,0.81}{#1}}
\newcommand{\BuiltInTok}[1]{#1}
\newcommand{\CharTok}[1]{\textcolor[rgb]{0.31,0.60,0.02}{#1}}
\newcommand{\CommentTok}[1]{\textcolor[rgb]{0.56,0.35,0.01}{\textit{#1}}}
\newcommand{\CommentVarTok}[1]{\textcolor[rgb]{0.56,0.35,0.01}{\textbf{\textit{#1}}}}
\newcommand{\ConstantTok}[1]{\textcolor[rgb]{0.56,0.35,0.01}{#1}}
\newcommand{\ControlFlowTok}[1]{\textcolor[rgb]{0.13,0.29,0.53}{\textbf{#1}}}
\newcommand{\DataTypeTok}[1]{\textcolor[rgb]{0.13,0.29,0.53}{#1}}
\newcommand{\DecValTok}[1]{\textcolor[rgb]{0.00,0.00,0.81}{#1}}
\newcommand{\DocumentationTok}[1]{\textcolor[rgb]{0.56,0.35,0.01}{\textbf{\textit{#1}}}}
\newcommand{\ErrorTok}[1]{\textcolor[rgb]{0.64,0.00,0.00}{\textbf{#1}}}
\newcommand{\ExtensionTok}[1]{#1}
\newcommand{\FloatTok}[1]{\textcolor[rgb]{0.00,0.00,0.81}{#1}}
\newcommand{\FunctionTok}[1]{\textcolor[rgb]{0.13,0.29,0.53}{\textbf{#1}}}
\newcommand{\ImportTok}[1]{#1}
\newcommand{\InformationTok}[1]{\textcolor[rgb]{0.56,0.35,0.01}{\textbf{\textit{#1}}}}
\newcommand{\KeywordTok}[1]{\textcolor[rgb]{0.13,0.29,0.53}{\textbf{#1}}}
\newcommand{\NormalTok}[1]{#1}
\newcommand{\OperatorTok}[1]{\textcolor[rgb]{0.81,0.36,0.00}{\textbf{#1}}}
\newcommand{\OtherTok}[1]{\textcolor[rgb]{0.56,0.35,0.01}{#1}}
\newcommand{\PreprocessorTok}[1]{\textcolor[rgb]{0.56,0.35,0.01}{\textit{#1}}}
\newcommand{\RegionMarkerTok}[1]{#1}
\newcommand{\SpecialCharTok}[1]{\textcolor[rgb]{0.81,0.36,0.00}{\textbf{#1}}}
\newcommand{\SpecialStringTok}[1]{\textcolor[rgb]{0.31,0.60,0.02}{#1}}
\newcommand{\StringTok}[1]{\textcolor[rgb]{0.31,0.60,0.02}{#1}}
\newcommand{\VariableTok}[1]{\textcolor[rgb]{0.00,0.00,0.00}{#1}}
\newcommand{\VerbatimStringTok}[1]{\textcolor[rgb]{0.31,0.60,0.02}{#1}}
\newcommand{\WarningTok}[1]{\textcolor[rgb]{0.56,0.35,0.01}{\textbf{\textit{#1}}}}
\usepackage{graphicx}
\makeatletter
\newsavebox\pandoc@box
\newcommand*\pandocbounded[1]{% scales image to fit in text height/width
  \sbox\pandoc@box{#1}%
  \Gscale@div\@tempa{\textheight}{\dimexpr\ht\pandoc@box+\dp\pandoc@box\relax}%
  \Gscale@div\@tempb{\linewidth}{\wd\pandoc@box}%
  \ifdim\@tempb\p@<\@tempa\p@\let\@tempa\@tempb\fi% select the smaller of both
  \ifdim\@tempa\p@<\p@\scalebox{\@tempa}{\usebox\pandoc@box}%
  \else\usebox{\pandoc@box}%
  \fi%
}
% Set default figure placement to htbp
\def\fps@figure{htbp}
\makeatother
\setlength{\emergencystretch}{3em} % prevent overfull lines
\providecommand{\tightlist}{%
  \setlength{\itemsep}{0pt}\setlength{\parskip}{0pt}}
\usepackage{amsmath}
% Shared LaTeX preamble for all Data Mining / ISLR chapter documents.
% Provides \makecoverpage{<assignment title>} — a standalone title page
% with fixed student identity, so every chapter/lab looks uniform.

\usepackage{graphicx}

\newcommand{\makecoverpage}[1]{%
\begin{titlepage}
\centering
\vspace*{2cm}

{\Large \textbf{Adventist University of Central Africa}}\\[0.4cm]
{\large Master's in Big Data Analytics}\\[2cm]

{\Huge \textbf{#1}}\\[0.5cm]
{\large \textit{Introduction to Statistical Learning with Applications in R}}\\[3cm]

\begin{tabular}{rl}
\textbf{Programme:} & MSc Big Data Analytics \\
\textbf{Course Title:} & Data Mining \\
\textbf{Course Code:} & MSDA 9124 \\
\textbf{Lecturer:} & Dr. Pacifique Nizeyimana \\
\textbf{Student Name:} & Wayne Rubangisa \\
\textbf{Student ID:} & 101220 \\
\end{tabular}

\vfill

{\large September 2026}

\end{titlepage}
}

\usepackage{bookmark}
\IfFileExists{xurl.sty}{\usepackage{xurl}}{} % add URL line breaks if available
\urlstyle{same}
\hypersetup{
  hidelinks,
  pdfcreator={LaTeX via pandoc}}

\author{}
\date{\vspace{-2.5em}}

\begin{document}

\makecoverpage{Chapter 5 --- Resampling Methods \\ Exercises 5 and 9}

\begin{Shaded}
\begin{Highlighting}[]
\FunctionTok{library}\NormalTok{(ISLR2)}
\FunctionTok{library}\NormalTok{(boot)}
\end{Highlighting}
\end{Shaded}

\section{\texorpdfstring{Exercise 5 --- Validation-Set Error for
Logistic Regression on
\texttt{Default}}{Exercise 5 --- Validation-Set Error for Logistic Regression on Default}}\label{exercise-5-validation-set-error-for-logistic-regression-on-default}

\subsection{(a) Fit the full-data logistic
regression}\label{a-fit-the-full-data-logistic-regression}

\begin{Shaded}
\begin{Highlighting}[]
\FunctionTok{set.seed}\NormalTok{(}\DecValTok{1}\NormalTok{)}
\NormalTok{glm.fit }\OtherTok{\textless{}{-}} \FunctionTok{glm}\NormalTok{(default }\SpecialCharTok{\textasciitilde{}}\NormalTok{ income }\SpecialCharTok{+}\NormalTok{ balance, }\AttributeTok{data =}\NormalTok{ Default, }\AttributeTok{family =}\NormalTok{ binomial)}
\FunctionTok{summary}\NormalTok{(glm.fit)}
\end{Highlighting}
\end{Shaded}

\begin{verbatim}
## 
## Call:
## glm(formula = default ~ income + balance, family = binomial, 
##     data = Default)
## 
## Coefficients:
##               Estimate Std. Error z value Pr(>|z|)    
## (Intercept) -1.154e+01  4.348e-01 -26.545  < 2e-16 ***
## income       2.081e-05  4.985e-06   4.174 2.99e-05 ***
## balance      5.647e-03  2.274e-04  24.836  < 2e-16 ***
## ---
## Signif. codes:  0 '***' 0.001 '**' 0.01 '*' 0.05 '.' 0.1 ' ' 1
## 
## (Dispersion parameter for binomial family taken to be 1)
## 
##     Null deviance: 2920.6  on 9999  degrees of freedom
## Residual deviance: 1579.0  on 9997  degrees of freedom
## AIC: 1585
## 
## Number of Fisher Scoring iterations: 8
\end{verbatim}

Both \texttt{income} and \texttt{balance} come out highly significant,
with \texttt{balance} the far stronger predictor by t-statistic
magnitude --- consistent with what Chapter 4 found for this same
relationship.

\subsection{(b) Validation-set test
error}\label{b-validation-set-test-error}

The validation-set approach is conceptually simple: split the data in
half, fit only on one half, and check performance on the half the model
never saw --- the resulting error rate is a genuine estimate of test
error, unlike the training error computed on data the model already fit.

\begin{Shaded}
\begin{Highlighting}[]
\FunctionTok{set.seed}\NormalTok{(}\DecValTok{1}\NormalTok{)}
\NormalTok{n }\OtherTok{\textless{}{-}} \FunctionTok{nrow}\NormalTok{(Default)}
\NormalTok{train }\OtherTok{\textless{}{-}} \FunctionTok{sample}\NormalTok{(n, n }\SpecialCharTok{/} \DecValTok{2}\NormalTok{)}

\NormalTok{glm.fit }\OtherTok{\textless{}{-}} \FunctionTok{glm}\NormalTok{(default }\SpecialCharTok{\textasciitilde{}}\NormalTok{ income }\SpecialCharTok{+}\NormalTok{ balance, }\AttributeTok{data =}\NormalTok{ Default,}
                \AttributeTok{family =}\NormalTok{ binomial, }\AttributeTok{subset =}\NormalTok{ train)}
\NormalTok{glm.probs }\OtherTok{\textless{}{-}} \FunctionTok{predict}\NormalTok{(glm.fit, Default[}\SpecialCharTok{{-}}\NormalTok{train, ], }\AttributeTok{type =} \StringTok{"response"}\NormalTok{)}
\NormalTok{glm.pred }\OtherTok{\textless{}{-}} \FunctionTok{ifelse}\NormalTok{(glm.probs }\SpecialCharTok{\textgreater{}} \FloatTok{0.5}\NormalTok{, }\StringTok{"Yes"}\NormalTok{, }\StringTok{"No"}\NormalTok{)}

\NormalTok{val.error }\OtherTok{\textless{}{-}} \FunctionTok{mean}\NormalTok{(glm.pred }\SpecialCharTok{!=}\NormalTok{ Default[}\SpecialCharTok{{-}}\NormalTok{train, ]}\SpecialCharTok{$}\NormalTok{default)}
\NormalTok{val.error}
\end{Highlighting}
\end{Shaded}

\begin{verbatim}
## [1] 0.0254
\end{verbatim}

A validation error around 2.5\% --- notably higher than the
essentially-zero training error a full-data fit could give the illusion
of, since 5,000 held-out observations never influenced the fitted
coefficients at all.

\subsection{(c) Repeating with three different
splits}\label{c-repeating-with-three-different-splits}

\begin{Shaded}
\begin{Highlighting}[]
\NormalTok{val\_error }\OtherTok{\textless{}{-}} \ControlFlowTok{function}\NormalTok{(seed) \{}
  \FunctionTok{set.seed}\NormalTok{(seed)}
\NormalTok{  train }\OtherTok{\textless{}{-}} \FunctionTok{sample}\NormalTok{(n, n }\SpecialCharTok{/} \DecValTok{2}\NormalTok{)}
\NormalTok{  fit }\OtherTok{\textless{}{-}} \FunctionTok{glm}\NormalTok{(default }\SpecialCharTok{\textasciitilde{}}\NormalTok{ income }\SpecialCharTok{+}\NormalTok{ balance, }\AttributeTok{data =}\NormalTok{ Default,}
             \AttributeTok{family =}\NormalTok{ binomial, }\AttributeTok{subset =}\NormalTok{ train)}
\NormalTok{  probs }\OtherTok{\textless{}{-}} \FunctionTok{predict}\NormalTok{(fit, Default[}\SpecialCharTok{{-}}\NormalTok{train, ], }\AttributeTok{type =} \StringTok{"response"}\NormalTok{)}
\NormalTok{  pred }\OtherTok{\textless{}{-}} \FunctionTok{ifelse}\NormalTok{(probs }\SpecialCharTok{\textgreater{}} \FloatTok{0.5}\NormalTok{, }\StringTok{"Yes"}\NormalTok{, }\StringTok{"No"}\NormalTok{)}
  \FunctionTok{mean}\NormalTok{(pred }\SpecialCharTok{!=}\NormalTok{ Default[}\SpecialCharTok{{-}}\NormalTok{train, ]}\SpecialCharTok{$}\NormalTok{default)}
\NormalTok{\}}

\NormalTok{errors }\OtherTok{\textless{}{-}} \FunctionTok{sapply}\NormalTok{(}\FunctionTok{c}\NormalTok{(}\DecValTok{1}\NormalTok{, }\DecValTok{2}\NormalTok{, }\DecValTok{3}\NormalTok{), val\_error)}
\NormalTok{errors}
\end{Highlighting}
\end{Shaded}

\begin{verbatim}
## [1] 0.0254 0.0238 0.0264
\end{verbatim}

The three error rates all land in a fairly tight band (roughly
2.4\%--2.6\%), but they're not identical --- which is the point. Every
one of these numbers is a valid estimate of the same underlying test
error, yet each depends on the arbitrary luck of which 5,000
observations happened to land in the training half. That split-to-split
variability is exactly the validation set approach's core weakness:
throwing away half the data for training, on a single random split,
means both the fitted model and the resulting error estimate are noisier
than they'd need to be with a method like \(k\)-fold cross-validation
that reuses the data more efficiently.

\subsection{\texorpdfstring{(d) Adding a \texttt{student} dummy
variable}{(d) Adding a student dummy variable}}\label{d-adding-a-student-dummy-variable}

\begin{Shaded}
\begin{Highlighting}[]
\FunctionTok{set.seed}\NormalTok{(}\DecValTok{1}\NormalTok{)}
\NormalTok{train }\OtherTok{\textless{}{-}} \FunctionTok{sample}\NormalTok{(n, n }\SpecialCharTok{/} \DecValTok{2}\NormalTok{)}
\NormalTok{glm.fit.student }\OtherTok{\textless{}{-}} \FunctionTok{glm}\NormalTok{(default }\SpecialCharTok{\textasciitilde{}}\NormalTok{ income }\SpecialCharTok{+}\NormalTok{ balance }\SpecialCharTok{+}\NormalTok{ student, }\AttributeTok{data =}\NormalTok{ Default,}
                        \AttributeTok{family =}\NormalTok{ binomial, }\AttributeTok{subset =}\NormalTok{ train)}
\NormalTok{glm.probs }\OtherTok{\textless{}{-}} \FunctionTok{predict}\NormalTok{(glm.fit.student, Default[}\SpecialCharTok{{-}}\NormalTok{train, ], }\AttributeTok{type =} \StringTok{"response"}\NormalTok{)}
\NormalTok{glm.pred }\OtherTok{\textless{}{-}} \FunctionTok{ifelse}\NormalTok{(glm.probs }\SpecialCharTok{\textgreater{}} \FloatTok{0.5}\NormalTok{, }\StringTok{"Yes"}\NormalTok{, }\StringTok{"No"}\NormalTok{)}

\NormalTok{val.error.student }\OtherTok{\textless{}{-}} \FunctionTok{mean}\NormalTok{(glm.pred }\SpecialCharTok{!=}\NormalTok{ Default[}\SpecialCharTok{{-}}\NormalTok{train, ]}\SpecialCharTok{$}\NormalTok{default)}
\NormalTok{val.error.student}
\end{Highlighting}
\end{Shaded}

\begin{verbatim}
## [1] 0.026
\end{verbatim}

With the same seed (and so the same train/test split) as (b), adding
\texttt{student} gives a validation error of 0.026, essentially
identical to --- in this run, marginally \emph{higher} than --- 0.0254
without it. So there's no evidence here that including \texttt{student}
improves predictive performance. This isn't surprising in hindsight:
\texttt{student} is strongly correlated with \texttt{balance} (students
tend to carry different balances than non-students), so much of whatever
information \texttt{student} carries about default risk is already
captured once \texttt{balance} is in the model --- adding it doesn't
give logistic regression much new signal to work with, only extra
estimation noise.

\newpage

\section{\texorpdfstring{Exercise 9 --- Bootstrap Estimates for the
\texttt{Boston}
Data}{Exercise 9 --- Bootstrap Estimates for the Boston Data}}\label{exercise-9-bootstrap-estimates-for-the-boston-data}

Every part below is about the same core idea: \texttt{medv} (or a
statistic of it) is being estimated from one particular sample of 506
census tracts, and we want to know how much that estimate would bounce
around if we happened to have drawn a different sample of the same size
from the same population. Textbook formulas exist for some statistics
(the sample mean); the bootstrap gives us an answer even for the ones
that don't (the median, a percentile) by treating our one sample as a
stand-in for the population and resampling from it directly.

\subsection{\texorpdfstring{(a) Estimate of the population mean,
\(\hat\mu\)}{(a) Estimate of the population mean, \textbackslash hat\textbackslash mu}}\label{a-estimate-of-the-population-mean-hatmu}

\begin{Shaded}
\begin{Highlighting}[]
\NormalTok{Boston2 }\OtherTok{\textless{}{-}}\NormalTok{ ISLR2}\SpecialCharTok{::}\NormalTok{Boston}
\NormalTok{mu.hat }\OtherTok{\textless{}{-}} \FunctionTok{mean}\NormalTok{(Boston2}\SpecialCharTok{$}\NormalTok{medv)}
\NormalTok{mu.hat}
\end{Highlighting}
\end{Shaded}

\begin{verbatim}
## [1] 22.53281
\end{verbatim}

\subsection{\texorpdfstring{(b) Standard error of \(\hat\mu\), by
formula}{(b) Standard error of \textbackslash hat\textbackslash mu, by formula}}\label{b-standard-error-of-hatmu-by-formula}

\[
\text{SE}(\hat\mu) = \frac{s}{\sqrt{n}}
\]

where \(s\) is the sample standard deviation of \texttt{medv} and
\(n = 506\).

\begin{Shaded}
\begin{Highlighting}[]
\NormalTok{se.formula }\OtherTok{\textless{}{-}} \FunctionTok{sd}\NormalTok{(Boston2}\SpecialCharTok{$}\NormalTok{medv) }\SpecialCharTok{/} \FunctionTok{sqrt}\NormalTok{(}\FunctionTok{nrow}\NormalTok{(Boston2))}
\NormalTok{se.formula}
\end{Highlighting}
\end{Shaded}

\begin{verbatim}
## [1] 0.4088611
\end{verbatim}

This measures how much \(\hat\mu\) would be expected to vary, on
average, across repeated samples of 506 tracts drawn from the same
population --- about \$409 (recall \texttt{medv} is in \$1000s), against
a mean around \$22,530.

\subsection{\texorpdfstring{(c) Standard error of \(\hat\mu\), by
bootstrap}{(c) Standard error of \textbackslash hat\textbackslash mu, by bootstrap}}\label{c-standard-error-of-hatmu-by-bootstrap}

\begin{Shaded}
\begin{Highlighting}[]
\FunctionTok{set.seed}\NormalTok{(}\DecValTok{1}\NormalTok{)}
\NormalTok{boot.fn.mean }\OtherTok{\textless{}{-}} \ControlFlowTok{function}\NormalTok{(data, index) }\FunctionTok{mean}\NormalTok{(data[index])}
\NormalTok{boot.mean }\OtherTok{\textless{}{-}} \FunctionTok{boot}\NormalTok{(Boston2}\SpecialCharTok{$}\NormalTok{medv, boot.fn.mean, }\AttributeTok{R =} \DecValTok{10000}\NormalTok{)}
\NormalTok{boot.mean}
\end{Highlighting}
\end{Shaded}

\begin{verbatim}
## 
## ORDINARY NONPARAMETRIC BOOTSTRAP
## 
## 
## Call:
## boot(data = Boston2$medv, statistic = boot.fn.mean, R = 10000)
## 
## 
## Bootstrap Statistics :
##     original       bias    std. error
## t1* 22.53281 -0.002350771   0.4069546
\end{verbatim}

The bootstrap standard error (0.407) comes out essentially identical to
the formula-based estimate in (b) (0.4089) --- reassuring, since for
something as well-behaved as a sample mean, the classical formula is
already exact asymptotically. The bootstrap isn't doing anything the
formula couldn't already do here; its real value shows up in (f) and
(h), for statistics that have no comparably simple formula at all.

\subsection{(d) 95\% confidence interval for the mean, via
bootstrap}\label{d-95-confidence-interval-for-the-mean-via-bootstrap}

Using the approximation given in the exercise,
\([\hat\mu - 2\,\text{SE}(\hat\mu),\ \hat\mu + 2\,\text{SE}(\hat\mu)]\):

\begin{Shaded}
\begin{Highlighting}[]
\NormalTok{se.boot }\OtherTok{\textless{}{-}} \FunctionTok{sd}\NormalTok{(boot.mean}\SpecialCharTok{$}\NormalTok{t)}
\NormalTok{ci.boot }\OtherTok{\textless{}{-}} \FunctionTok{c}\NormalTok{(mu.hat }\SpecialCharTok{{-}} \DecValTok{2} \SpecialCharTok{*}\NormalTok{ se.boot, mu.hat }\SpecialCharTok{+} \DecValTok{2} \SpecialCharTok{*}\NormalTok{ se.boot)}
\NormalTok{ci.boot}
\end{Highlighting}
\end{Shaded}

\begin{verbatim}
## [1] 21.71890 23.34672
\end{verbatim}

\begin{Shaded}
\begin{Highlighting}[]
\FunctionTok{t.test}\NormalTok{(Boston2}\SpecialCharTok{$}\NormalTok{medv)}
\end{Highlighting}
\end{Shaded}

\begin{verbatim}
## 
##  One Sample t-test
## 
## data:  Boston2$medv
## t = 55.111, df = 505, p-value < 2.2e-16
## alternative hypothesis: true mean is not equal to 0
## 95 percent confidence interval:
##  21.72953 23.33608
## sample estimates:
## mean of x 
##  22.53281
\end{verbatim}

The bootstrap-based interval (roughly 21.72 to 23.35) is nearly
indistinguishable from \texttt{t.test()}'s interval --- expected, since
both are really estimating the same quantity (\(\pm 2\) standard errors
around the mean is a close approximation to the \(t\)-distribution
multiplier \texttt{t.test()} uses at \(n=506\)), just arrived at through
different machinery: one assumes normal-theory sampling distribution
math, the other estimates the same spread empirically by actually
resampling.

\subsection{(e) Estimate of the population
median}\label{e-estimate-of-the-population-median}

\begin{Shaded}
\begin{Highlighting}[]
\NormalTok{med.hat }\OtherTok{\textless{}{-}} \FunctionTok{median}\NormalTok{(Boston2}\SpecialCharTok{$}\NormalTok{medv)}
\NormalTok{med.hat}
\end{Highlighting}
\end{Shaded}

\begin{verbatim}
## [1] 21.2
\end{verbatim}

\subsection{(f) Standard error of the median, via
bootstrap}\label{f-standard-error-of-the-median-via-bootstrap}

There is no simple closed-form formula for the standard error of a
sample median (unlike the mean, in (b)) --- this is precisely the
situation the bootstrap is built for.

\begin{Shaded}
\begin{Highlighting}[]
\FunctionTok{set.seed}\NormalTok{(}\DecValTok{1}\NormalTok{)}
\NormalTok{boot.fn.median }\OtherTok{\textless{}{-}} \ControlFlowTok{function}\NormalTok{(data, index) }\FunctionTok{median}\NormalTok{(data[index])}
\NormalTok{boot.median }\OtherTok{\textless{}{-}} \FunctionTok{boot}\NormalTok{(Boston2}\SpecialCharTok{$}\NormalTok{medv, boot.fn.median, }\AttributeTok{R =} \DecValTok{10000}\NormalTok{)}
\NormalTok{boot.median}
\end{Highlighting}
\end{Shaded}

\begin{verbatim}
## 
## ORDINARY NONPARAMETRIC BOOTSTRAP
## 
## 
## Call:
## boot(data = Boston2$medv, statistic = boot.fn.median, R = 10000)
## 
## 
## Bootstrap Statistics :
##     original    bias    std. error
## t1*     21.2 -0.014705   0.3780355
\end{verbatim}

The estimated standard error (0.378) is somewhat \emph{smaller} than the
mean's standard error from (c) --- the median is less sensitive than the
mean to the handful of very high-\texttt{medv} tracts (recall
\texttt{medv} appears capped at \$50,000 in this data), so it varies a
bit less across resamples.

\subsection{(g) Estimate of the 10th
percentile}\label{g-estimate-of-the-10th-percentile}

\begin{Shaded}
\begin{Highlighting}[]
\NormalTok{mu.hat}\FloatTok{.01} \OtherTok{\textless{}{-}} \FunctionTok{quantile}\NormalTok{(Boston2}\SpecialCharTok{$}\NormalTok{medv, }\FloatTok{0.1}\NormalTok{)}
\NormalTok{mu.hat}\FloatTok{.01}
\end{Highlighting}
\end{Shaded}

\begin{verbatim}
##   10% 
## 12.75
\end{verbatim}

\subsection{(h) Standard error of the 10th percentile, via
bootstrap}\label{h-standard-error-of-the-10th-percentile-via-bootstrap}

\begin{Shaded}
\begin{Highlighting}[]
\FunctionTok{set.seed}\NormalTok{(}\DecValTok{1}\NormalTok{)}
\NormalTok{boot.fn.q10 }\OtherTok{\textless{}{-}} \ControlFlowTok{function}\NormalTok{(data, index) }\FunctionTok{quantile}\NormalTok{(data[index], }\FloatTok{0.1}\NormalTok{)}
\NormalTok{boot.q10 }\OtherTok{\textless{}{-}} \FunctionTok{boot}\NormalTok{(Boston2}\SpecialCharTok{$}\NormalTok{medv, boot.fn.q10, }\AttributeTok{R =} \DecValTok{10000}\NormalTok{)}
\NormalTok{boot.q10}
\end{Highlighting}
\end{Shaded}

\begin{verbatim}
## 
## ORDINARY NONPARAMETRIC BOOTSTRAP
## 
## 
## Call:
## boot(data = Boston2$medv, statistic = boot.fn.q10, R = 10000)
## 
## 
## Bootstrap Statistics :
##     original   bias    std. error
## t1*    12.75 0.007625   0.5014507
\end{verbatim}

The standard error here (0.5015) is the largest of the three statistics
examined --- larger than both the mean's and the median's. This tracks
with intuition: the 10th percentile is estimated from comparatively few,
more extreme observations in the lower tail of \texttt{medv}, so there's
inherently less data locally supporting that estimate than there is
supporting the mean or median, which draw on the full distribution. This
is exactly the kind of finding the bootstrap makes practical to obtain
--- no standard formula for the standard error of a percentile exists at
all, yet the same resampling procedure that worked for the mean and
median hands us an answer here just as directly.

\end{document}
